% =========================================================================
% TRIP4/ASCC Proteomics Poster — LaTeX Template
% =========================================================================
% This template ensures UNIFORM FONT SIZES across all figures.
%
% HOW IT WORKS:
% 1. R generates figures as PDF (vector format) via save_poster_figure()
% 2. LaTeX includes them via \includegraphics
% 3. When LaTeX resizes a figure, the TEXT INSIDE the figure scales
%    proportionally with the figure — so relative font sizes stay consistent
% 4. For TRUE pixel-identical fonts: generate figures at the EXACT size
%    they appear in the poster (see width values below)
%
% REQUIREMENTS:
%   - TeX Live or MiKTeX with beamerposter package
%   - Install: tlmgr install beamerposter (TeX Live)
%   Or on Windows MiKTeX: it auto-installs on first compile
%
% USAGE:
%   1. Run R pipeline: make aruna-all (generates PDFs in output/figures/)
%   2. Copy desired PDFs to poster/figures/ (or update paths below)
%   3. Compile: pdflatex poster.tex (run 2-3x for layout)
%
% ALTERNATIVE (no beamerposter):
%   Use the standalone article-based version (poster_simple.tex)
% =========================================================================

\documentclass[final]{beamer}
\usepackage[size=custom,width=120,height=90,scale=1.0]{beamerposter}
% 120cm x 90cm = standard portrait poster size
% Change to custom,width=180,height=100 for landscape

\usepackage[utf8]{inputenc}
\usepackage[T1]{fontenc}
\usepackage{lmodern}              % Clean Latin Modern fonts
\usepackage{graphicx}             % Include figures
\usepackage{booktabs}             % Nice tables
\usepackage{tikz}                 % Diagrams
\usepackage{enumitem}             % List formatting
\usepackage{xcolor}

% ---- Color scheme (match R GLOBAL_COLORS) ----
\definecolor{asccblue}{HTML}{0072B2}
\definecolor{interactorgreen}{HTML}{009E73}
\definecolor{enrichedred}{HTML}{D55E00}
\definecolor{bglight}{HTML}{F5F5F5}

% ---- Beamer theme ----
\usetheme{default}
\setbeamertemplate{navigation symbols}{}  % Remove nav arrows
\setbeamercolor{block title}{fg=white, bg=asccblue}
\setbeamercolor{block body}{fg=black, bg=bglight}
\setbeamerfont{block title}{size=\large, series=\bfseries}

% ---- Block margins ----
\setbeamertemplate{block begin}{
  \vskip0.5em
  \begin{beamercolorbox}[ht=2.8em, leftskip=1em, colsep*=.75ex]{block title}
    \usebeamerfont{block title}\insertblocktitle
  \end{beamercolorbox}
  {\parskip0pt\par}
  \ifbeamercolorempty[bg]{block body}{}{\nointerlineskip\vskip-0.5pt}
  \usebeamerfont{block body}
  \begin{beamercolorbox}[colsep*=.75ex, sep=.75ex, vmode]{block body}
}
\setbeamertemplate{block end}{
  \end{beamercolorbox}
  \vskip0.5em
}

% =========================================================================
% POSTER CONTENT — EDIT BELOW
% =========================================================================

\title{TRIP4/ASCC Complex: Proximity Labeling and Interaction Proteomics}
\author{A.\,S.\,[Your Name] \quad F.\,[Collaborators]}
\institute{[Your Institute]}
\date{}

\begin{document}

\begin{frame}[t]

% ---- Title bar ----
\vskip1cm
\begin{center}
  {\veryHuge\bfseries\color{asccblue} TRIP4/ASCC Complex: Proximity Labeling and Interaction Proteomics\par}
  \vskip0.5cm
  {\Large [Your Name] $\cdot$ [Institute] $\cdot$ [Year]\par}
\end{center}
\vskip1cm

% ---- Two-column layout ----
\begin{columns}[t]

% ==================== LEFT COLUMN ====================
\begin{column}{0.48\textwidth}

\begin{block}{Introduction}
\small
TRIP4 (Transcription Intermediary Factor 4) is the core subunit of the ASCC
complex, involved in RNA processing and stress response. We use TurboID
proximity labeling and Flag IP to map the TRIP4 interaction network.
\end{block}

\begin{block}{Results 1: TurboID Volcano}
% Include figure as PDF (vector). Width controls scaling.
% Font sizes inside figure scale proportionally — consistent across all panels.
\centering
\includegraphics[width=0.95\linewidth]{figures/lydia_volcano_network.pdf}
\end{block}

\begin{block}{Results 2: GO Enrichment (Biological Process)}
\centering
\includegraphics[width=0.95\linewidth]{figures/targeted_GO_turbo_BP_dotplot.pdf}
\end{block}

\begin{block}{Results 3: STRING Network}
\centering
\includegraphics[width=0.95\linewidth]{figures/ra_string_style_network.pdf}
\end{block}

\end{column}

% ==================== RIGHT COLUMN ====================
\begin{column}{0.48\textwidth}

\begin{block}{Flag IP Validation}
% NOTE: Both figures use SAME width (0.95\linewidth).
% This ensures fonts are the SAME SIZE in both panels.
\centering
\includegraphics[width=0.95\linewidth]{figures/flagip_volcano_overlap.pdf}
\vskip0.3cm
\includegraphics[width=0.95\linewidth]{figures/targeted_venn_turbo_vs_flag.pdf}
\end{block}

\begin{block}{KEGG Pathway Enrichment}
\centering
\includegraphics[width=0.95\linewidth]{figures/targeted_KEGG_turbo_dotplot.pdf}
\end{block}

\begin{block}{Validated Interactors (C-Flag + N-Flag)}
\centering
\includegraphics[width=0.95\linewidth]{figures/flagip_GO_validated_both_BP_dotplot.pdf}
\end{block}

\begin{block}{Conclusions}
\small
\begin{itemize}[leftmargin=1em, itemsep=0.3em]
  \item TRIP4 proximity labeling identifies ASCC complex and RNA-binding proteins
  \item Flag IP (C-Flag and N-Flag) validates key interactors
  \item GO/KEGG enrichment reveals RNA processing pathways
  \item Retinoic acid treatment reshapes the interaction network
\end{itemize}
\end{block}

\end{column}
\end{columns}

\end{frame}

\end{document}
